\documentclass[11pt]{article}
\usepackage[margin=1in]{geometry}
\usepackage{amsmath,amssymb,amsfonts}
\usepackage{bm}
\usepackage{booktabs}
\usepackage{enumitem}
\usepackage{xcolor}
\usepackage{hyperref}
\hypersetup{colorlinks=true,linkcolor=blue!60!black,citecolor=blue!60!black,urlcolor=blue!60!black}

\newcommand{\pderiv}[2]{\frac{\partial #1}{\partial #2}}
\newcommand{\Rav}{\widetilde{\mathrm{Ra}}}
\newcommand{\avg}[1]{\langle #1 \rangle}
\newcommand{\kh}{|\bm{k}|}
\newcommand{\Dzone}{D^{(1)}_z}
\newcommand{\Dztwo}{D^{(2)}_z}
\newcommand{\Hmat}{H}
\newcommand{\Emat}{E_D}
\newcommand{\Pmat}{P_N}
\newcommand{\Gmat}{G}

\title{\textbf{Vertical Operator Experiments for the FD NHQG Benchmark}\\[8pt]
\large Compact Schemes, SBP Structure, Compatibility, and Nonnormality}
\author{}
\date{April 2026}

\begin{document}
\maketitle
\vspace{0.5em}

\section{Purpose}

This note records the operator-level reasoning behind the finite-difference-in-$z$
benchmark branch in \texttt{fd\_vertical\_benchmark/}, with particular focus on
two replacement families for the legacy centered second-order vertical
operators:
\begin{itemize}[nosep]
\item compact ``spectral-tending'' operators,
\item summation-by-parts (SBP) operators designed around a discrete energy law.
\end{itemize}

The key lesson from the first compact experiment is subtle:
\begin{itemize}[nosep]
\item a compact first derivative can be \emph{much more accurate} as a local
  derivative approximation,
\item yet still be a \emph{worse} building block for the NHQG coupled implicit
  operator if it is not paired with a compatible second derivative and
  compatible boundary closures.
\end{itemize}

So the goal is not simply to make the stencil higher-order. The goal is to make
the \emph{semi-discrete operator family} coherent.

\section{Where the Vertical Operators Enter}

In the FD benchmark, the prognostic fluctuation fields
$\psi(z,\bm{k},t)$, $w(z,\bm{k},t)$, $\theta(z,\bm{k},t)$ are stored on a
uniform vertical grid. The mean temperature deviation $\Theta'_b(z,t)$ is also
stored on the same vertical grid, but only at interior Dirichlet nodes.

The relevant discrete maps are:
\begin{itemize}[nosep]
\item $\Emat$: interior Dirichlet extension, inserting zero boundary values,
\item $\Pmat$: Neumann reconstruction for $\psi$, recovering boundary values
  from interior values so that the discrete derivative vanishes at $z=0,1$,
\item $\Dzone$: full-grid first-derivative matrix,
\item $\Dztwo$: interior Dirichlet second-derivative matrix.
\end{itemize}

From these, the solver constructs
\begin{align}
D^{(\psi)}_z &= S_{\mathrm{int}}\,\Dzone\,\Pmat,
\\
D^{(D)}_z &= S_{\mathrm{int}}\,\Dzone\,\Emat,
\end{align}
where $S_{\mathrm{int}}$ restricts a full-grid quantity to the interior nodes.

The fluctuation coupling that matters most for the IMEX step is
\begin{equation}
\boxed{
B \;=\; D^{(\psi)}_z D^{(D)}_z.
}
\label{eq:B_def}
\end{equation}

For each horizontal shell $\kh^2$, the FD benchmark precomputes the dense
interior solve
\begin{equation}
\boxed{
A(\kh) \;=\; \alpha_w^{\mathrm{eff}} I
- (\gamma \Delta t)^2 \frac{1}{\kh^2\,\alpha_q}\, B.
}
\label{eq:A_shell}
\end{equation}

The mean-temperature solve uses a Dirichlet second derivative:
\begin{equation}
\boxed{
\left(I - \gamma \Delta t\, \frac{\epsilon_b^2}{\sigma}\, \Dztwo \right)\Theta_b'^{\,*}
= \text{RHS}.
}
\label{eq:mean_temp_solve}
\end{equation}

Therefore the vertical operators matter in \emph{three} logically distinct ways:
\begin{enumerate}[nosep]
\item local derivative accuracy,
\item structural properties of the coupled operator $B$,
\item dissipativity of the mean-temperature diffusion operator $\Dztwo$.
\end{enumerate}

\section{What Properties We Actually Need}

\subsection{Consistency and formal order}

At minimum, $\Dzone$ and $\Dztwo$ should approximate $\partial_z$ and
$\partial_{zz}$ with the intended order on smooth functions. This matters,
especially for a benchmark branch whose purpose is to isolate vertical
discretization effects from the horizontal pseudospectral part.

But formal order alone is not enough.

\subsection{Boundary-condition compatibility}

The FD benchmark does not evolve full-grid boundary values independently.
Instead:
\begin{itemize}[nosep]
\item $w$, $\theta$, and $\Theta_b'$ use Dirichlet extension,
\item $\psi$ uses a reconstructed Neumann boundary map.
\end{itemize}

So the boundary treatment is not an add-on. It is part of the operator.
If $\Dzone$ is changed but $\Pmat$ and $\Dztwo$ remain tied to an incompatible lower
order discretization, the resulting coupled matrices can become highly
nonnormal even when each separate piece looks reasonable in isolation.

\subsection{Projected versus direct Neumann reduction}

For a compact first derivative, it is often cleaner to think in terms of the
implicit relation
\begin{equation}
\boxed{
A^{(1)} d = B^{(1)} f,
}
\end{equation}
where $f$ is the full-grid field and $d$ is the full-grid discrete derivative.

In the original FD benchmark implementation, the workflow was:
\begin{enumerate}[nosep]
\item solve for the full operator $\Dzone = (A^{(1)})^{-1} B^{(1)}$,
\item reconstruct boundary values through a separate Neumann map $\Pmat$,
\item then form the reduced derivative
  $D_z^{(\psi)} = S_{\mathrm{int}} \Dzone \Pmat$.
\end{enumerate}

That is algebraically valid, but it separates the compact inversion from the
boundary elimination.

The more direct construction treats the interior field values $f_i$ as the
known data, imposes the Neumann conditions $d_b = 0$ at the two boundary nodes,
and solves in one shot for the unknown interior derivatives $d_i$ and boundary
values $f_b$. Partitioning the compact relation into boundary and interior
blocks gives
\begin{equation}
\begin{bmatrix}
A^{(1)}_{bi} & -B^{(1)}_{bb} \\
A^{(1)}_{ii} & -B^{(1)}_{ib}
\end{bmatrix}
\begin{bmatrix}
d_i \\ f_b
\end{bmatrix}
=
\begin{bmatrix}
B^{(1)}_{bi} \\ B^{(1)}_{ii}
\end{bmatrix}
f_i.
\label{eq:direct_neumann_reduced}
\end{equation}

Solving \eqref{eq:direct_neumann_reduced} directly produces both the reduced
Neumann derivative $D_z^{(\psi)}$ and the reconstruction map $\Pmat$ from the
\emph{same} linear system. This is conceptually cleaner because the boundary
conditions are built into the reduced operator itself.

For the present fourth-order compact stencil, the direct and projected
constructions turned out to be the same operator to numerical precision:
\begin{itemize}[nosep]
\item ${\Pmat}_{\mathrm{direct}} - {\Pmat}_{\mathrm{projected}} = 0$ to machine
  precision in the current audit,
\item $D_{z,\mathrm{direct}}^{(\psi)}$ and
  $D_{z,\mathrm{projected}}^{(\psi)}$ differ at only roundoff level.
\end{itemize}

So the direct construction is best understood as a \emph{cleaner formulation},
not as a new discretization. If the compact branch remains unstable after this
rewrite, that points away from bookkeeping error and back toward the operator
family itself, especially the nonnormality already diagnosed in $B$.

\subsection{A dissipative second derivative}

For a Dirichlet diffusion operator, the fundamental discrete property is not
just spectral negativity. It is a discrete energy inequality.

Let $\Hmat_{\mathrm{int}}$ be the interior mass matrix and $\Hmat_f$ the full-grid
quadrature matrix. If $\Gmat$ is the discrete gradient used for Dirichlet
fields, then the best possible property is
\begin{equation}
\boxed{
D^{(\mathrm{energy})}_{zz}
\,=\,
-\,\Hmat_{\mathrm{int}}^{-1}\Gmat^T \Hmat_f \Gmat.
}
\label{eq:D2_energy}
\end{equation}

Then for any interior vector $u$,
\begin{equation}
u^T \Hmat_{\mathrm{int}} D^{(\mathrm{energy})}_{zz} u
=
- (\Gmat u)^T \Hmat_f (\Gmat u)
\le 0.
\label{eq:energy_dissipation}
\end{equation}

This is the discrete analogue of
\(
\int_0^1 u\,u_{zz}\,dz = -\int_0^1 |u_z|^2\,dz
\)
for homogeneous Dirichlet data.

For the present FD benchmark, however, there is an important modeling priority:
\(\Dztwo\) enters only through the implicit mean-temperature equation, while the
observed instability is primarily a fluctuation problem tied to the first
derivative coupling \( B = D_z^{(\psi)} D_z^{(D)} \). Therefore the main
experimental priority is:
\begin{itemize}[nosep]
\item invest operator sophistication in \(\Dzone\),
\item keep \(\Dztwo\) simple and robust unless there is evidence that the
  mean-temperature diffusion itself is the culprit.
\end{itemize}

That is why, after the initial SBP implementation, the main comparison branch
is taken to be
\[
\boxed{
\texttt{sbp42 + centered2}
}
\]
rather than \(\texttt{sbp42 + sbp42\_energy}\). The wide SBP-induced
\(\Dztwo\) remains useful as a diagnostic branch, but not as the main
production candidate.

\subsection{Why an SBP branch is the natural fallback}

The compact branch taught us that local derivative accuracy is not enough.
What we really need is a first derivative together with boundary closures and a
second derivative that cooperate in an energy estimate.

This is exactly what SBP operators are built to provide. A first-derivative SBP
operator consists of a positive-definite norm matrix $\Hmat$ and a derivative
matrix $\Dzone$ such that
\begin{equation}
\boxed{
\Hmat \Dzone + {\Dzone}^T \Hmat = B_{\partial},
}
\label{eq:sbp_identity}
\end{equation}
where
\(
B_{\partial} = \mathrm{diag}(-1,0,\dots,0,1).
\)

Equation \eqref{eq:sbp_identity} is the discrete analogue of integration by
parts. For the FD NHQG benchmark this matters in two ways:
\begin{enumerate}[nosep]
\item the boundary contribution is explicit and easier to audit than the
  compact branch's amplification through $\Pmat$,
\item the same norm matrix $\Hmat$ can be used to build a dissipative Dirichlet
  Laplacian through \eqref{eq:D2_energy}.
\end{enumerate}

So the SBP branch is not trying to be more ``spectral'' than compact Pad\'e.
It is trying to be more \emph{structurally trustworthy}.

\subsection{The concrete SBP operator used here}

The first SBP implementation in the benchmark uses the standard diagonal-norm
fourth-order interior, second-order boundary operator, often denoted
$D_1(4,2)$ or ``SBP42''.

Its norm matrix is
\begin{equation}
\boxed{
\Hmat = h\,\mathrm{diag}\!\left(
\frac{17}{48},\,
\frac{59}{48},\,
\frac{43}{48},\,
\frac{49}{48},\,
1,\dots,1,\,
\frac{49}{48},\,
\frac{43}{48},\,
\frac{59}{48},\,
\frac{17}{48}
\right).
}
\label{eq:sbp42_H}
\end{equation}

The interior stencil is the classical fourth-order centered derivative,
\begin{equation}
\boxed{
(\Dzone f)_i
=
\frac{1}{h}\left(
\frac{1}{12} f_{i-2}
- \frac{2}{3} f_{i-1}
+ \frac{2}{3} f_{i+1}
- \frac{1}{12} f_{i+2}
\right),
\qquad 4 \le i \le N-4.
}
\label{eq:sbp42_interior}
\end{equation}

The leading boundary rows are
\begin{align}
(\Dzone f)_0 &=
\frac{1}{h}\left(
-\frac{24}{17} f_0
+\frac{59}{34} f_1
-\frac{4}{17} f_2
-\frac{3}{34} f_3
\right),
\\
(\Dzone f)_1 &=
\frac{1}{h}\left(
-\frac{1}{2} f_0
+\frac{1}{2} f_2
\right),
\\
(\Dzone f)_2 &=
\frac{1}{h}\left(
\frac{4}{43} f_0
-\frac{59}{86} f_1
+\frac{59}{86} f_3
-\frac{4}{43} f_4
\right),
\\
(\Dzone f)_3 &=
\frac{1}{h}\left(
\frac{3}{98} f_0
-\frac{59}{98} f_2
+\frac{32}{49} f_4
-\frac{4}{49} f_5
\right),
\end{align}
with mirrored rows at the upper boundary. These coefficients satisfy the SBP
identity \eqref{eq:sbp_identity} exactly.

In the current benchmark, the SBP second derivative is not a separate narrow
operator with its own optimized closure. Instead, it is the compatible
energy-form Laplacian induced by the same SBP gradient:
\begin{equation}
\boxed{
{\Dztwo}^{(\mathrm{sbp})}
=
-\Hmat_{\mathrm{int}}^{-1}\Gmat^T \Hmat_f \Gmat,
\qquad
\Gmat = \Dzone \Emat.
}
\label{eq:sbp42_d2}
\end{equation}

This is a \emph{wide} operator rather than a narrow optimized SBP second
derivative. Its value is that it is guaranteed dissipative in the SBP norm,
even if its raw Laplacian accuracy is not spectacular. That is acceptable for a
first SBP branch whose job is to test whether the instability is tied more to
boundary-transient structure than to local stencil accuracy.

\subsection{Controlled nonnormality}

The first compact experiment showed why eigenvalues alone are not enough.
Even if $\mathrm{spec}(B)$ remains entirely in the left half-plane, the solver
can still suffer large transient amplification if $B$ is strongly nonnormal.

The most useful operator-level indicators are:
\begin{itemize}[nosep]
\item the commutator norm
  $\|B^T B - B B^T\|$,
\item the eigenvector condition number $\kappa(V)$ for $B = V \Lambda V^{-1}$,
\item the largest eigenvalue of the symmetric part
  $\tfrac{1}{2}(B + B^T)$,
\item the induced amplification in the Neumann reconstruction $\Pmat$.
\end{itemize}

These quantify how much a stable spectrum can still support large transient
growth.

\subsection{Good shellwise IMEX conditioning}

Even if $B$ is ugly, the actual implicit stage can still be fine if $A(\kh)$ in
\eqref{eq:A_shell} stays well-conditioned. This distinction turned out to be
important in the audit below.

\section{What the First Compact-$D^{(1)}_z$ Experiment Showed}

The first compact experiment replaced the legacy centered second-order
$\Dzone$ with a fourth-order tridiagonal compact first derivative, while
leaving the mean-temperature $\Dztwo$ on the old second-order Dirichlet
Laplacian.

\subsection{What improved}

As a plain derivative approximation, the compact operator behaved exactly as one
would hope:
\begin{itemize}[nosep]
\item for smooth trigonometric probes, the $\partial_z$ error dropped by orders
  of magnitude relative to the centered second-order stencil,
\item the reconstructed Neumann derivative for a smooth cosine profile was also
  dramatically more accurate.
\end{itemize}

So if one only looks at derivative accuracy, the compact operator appears to be
an immediate win.

\subsection{What became worse}

The operator audit at $N_z=128$ told a different story. The compact
$\Dzone$ branch had:
\begin{center}
\begin{tabular}{lcc}
\toprule
Metric & centered2 & compact4 \\
\midrule
$\|\Dzone\|_{\infty}$ & $5.12\times 10^2$ & $1.74\times 10^3$ \\
$\|\Pmat\|_{\infty}$ & $1.67$ & $5.32$ \\
boundary row $\ell^1$ norm & $1.67$ & $5.32$ \\
$\|B^T B - B B^T\|_2$ & $5.52\times 10^7$ & $4.83\times 10^9$ \\
$\kappa(V_B)$ & $2.52$ & $2.19\times 10^2$ \\
$\lambda_{\max}\!\left(\tfrac{B+B^T}{2}\right)$ & $2.68\times 10^2$ & $2.42\times 10^4$ \\
\bottomrule
\end{tabular}
\end{center}

At the same time:
\begin{itemize}[nosep]
\item $\mathrm{Re}\,\lambda(B)$ stayed non-positive to numerical precision,
\item the shellwise IMEX matrices $A(\kh)$ stayed extremely well-conditioned
  (condition numbers essentially $1$).
\end{itemize}

This is a very specific diagnosis:
\begin{equation}
\boxed{
\text{the failure was not caused by a singular IMEX solve;}
\quad
\text{it was caused by a much more nonnormal and boundary-amplifying } B.
}
\end{equation}

\subsection{Observed run-level consequence}

The FD benchmark with the old centered second-order branch remained finite
through $t=3.75$ and first produced saved NaNs at $t=4.0$.
The compact-$\Dzone$ branch, with the old $\Dztwo$ left untouched, produced NaNs by
$t=2.5$.

So the first compact attempt was a clean negative result, but it was also
informative. It strongly suggests that the problem is \emph{compatibility of
the operator family}, not the mere use of compact stencils.

\section{How to Build a Compatible $D^{(2)}_z$}

The next step is to add a second derivative that belongs to the same vertical
operator family as the compact first derivative, instead of mixing compact
$\Dzone$ with the old second-order $\Dztwo$.

\subsection{Raw compact second derivative}

The direct analogue of the compact first derivative is a tridiagonal compact
second derivative on the full grid:
\begin{equation}
A_2 f'' = B_2 f,
\qquad
\boxed{
D^{(\mathrm{raw})}_{zz,f} = A_2^{-1} B_2.
}
\label{eq:D2_raw_full}
\end{equation}

For the interior stencil we use the standard fourth-order compact relation
\begin{equation}
\frac{1}{10} f''_{j-1} + f''_j + \frac{1}{10} f''_{j+1}
=
\frac{6}{5}\,\frac{f_{j-1} - 2f_j + f_{j+1}}{h^2}.
\label{eq:compact_second_interior}
\end{equation}

At the boundary, a one-sided closure is constructed by matching polynomials
through degree $5$:
\begin{equation}
f''_0 + 10 f''_1
=
\frac{1}{h^2}
\left(
\frac{145}{12} f_0
- \frac{76}{3} f_1
+ \frac{29}{2} f_2
- \frac{4}{3} f_3
+ \frac{1}{12} f_4
\right),
\label{eq:compact_second_boundary}
\end{equation}
with the top closure obtained by reflection.

For Dirichlet fields, the interior operator is then restricted by zero
extension:
\begin{equation}
\boxed{
D^{(\mathrm{raw})}_{zz,\mathrm{dir}}
=
S_{\mathrm{int}}\,D^{(\mathrm{raw})}_{zz,f}\,\Emat.
}
\label{eq:D2_raw_dir}
\end{equation}

This branch is valuable because it tells us what the ``classical'' compact
second derivative does on its own.

\subsection{Energy-compatible second derivative}

The safer branch is the compatibility construction already motivated in
\eqref{eq:D2_energy}.

Let
\begin{equation}
\boxed{
\Gmat = D^{(1)}_{z,f} \Emat
}
\label{eq:G_def}
\end{equation}
be the full-grid gradient of a Dirichlet interior field, using the \emph{same}
full-grid first derivative $D^{(1)}_{z,f}$ as the rest of the compact family.

Let $\Hmat_f$ be the full-grid trapezoidal mass matrix and
$\Hmat_{\mathrm{int}}$ the interior mass matrix. Then define
\begin{equation}
\boxed{
D^{(\mathrm{energy})}_{zz,\mathrm{dir}}
=
- \Hmat_{\mathrm{int}}^{-1}\Gmat^T \Hmat_f \Gmat
}
\label{eq:D2_compatible}
\end{equation}

This is not ``raw compact'' in the Lele sense. It is better viewed as the
Dirichlet Laplacian \emph{induced by the chosen first derivative and discrete
inner product}. Its decisive advantage is the exact dissipation identity
\eqref{eq:energy_dissipation}.

\subsection{Why both branches are worth keeping}

The two branches answer different questions:
\begin{itemize}[nosep]
\item \texttt{compact4\_raw}: does a classical compact-$D^{(2)}$ closure fix the
  mixed-operator pathology?
\item \texttt{from\_d1\_energy}: if not, does a dissipative operator induced by
  the same compact gradient cure the incompatibility?
\end{itemize}

That is a better experiment than jumping directly from the failed first attempt
to a single new ``improved'' operator with no control comparison.

\section{Audit Checklist for the New $D^{(2)}_z$}

Any candidate $\Dztwo$ should pass the following operator-level checks before a
long NHQG run:
\begin{enumerate}[nosep]
\item $\Dztwo$ is consistent on smooth probes such as $\sin(\pi z)$.
\item $\mathrm{spec}(\Dztwo)$ has non-positive real part.
\item $I - \alpha \Dztwo$ is well-conditioned for the relevant implicit
  coefficients.
\item The energy form $u^T \Hmat_{\mathrm{int}} \Dztwo u$ is non-positive for the
  energy-compatible branch, and at least empirically non-positive on random
  vectors for the raw branch.
\item The mixed operator $B = D^{(\psi)}_z D^{(D)}_z$ does not become more
  nonnormal after the full operator family is assembled.
\item The benchmark run no longer exhibits an earlier blow-up than the legacy
  centered branch.
\end{enumerate}

\section{Bottom Line}

The first compact experiment failed for a useful reason:
\begin{equation}
\boxed{
\text{higher local accuracy of } \Dzone
\;\not\Rightarrow\;
\text{better semi-discrete NHQG dynamics.}
}
\end{equation}

The properties that matter are:
\begin{itemize}[nosep]
\item boundary compatibility,
\item dissipativity of $\Dztwo$,
\item moderate nonnormality of the coupled operator,
\item good shellwise conditioning of the IMEX stage solve.
\end{itemize}

This leads directly to the next implementation strategy:
\begin{enumerate}[nosep]
\item keep the compact first derivative,
\item add both a raw compact and an energy-compatible second derivative,
\item audit both at the operator level,
\item only then return to long NHQG benchmark runs.
\end{enumerate}

\end{document}
